10.5 Higher-Order Derivatives
81
If the mapping (10.58) is differentiable in $U$, then the derivative of $f$, given by
$$f': U \to \mathcal{L}(X; Y),$$
(10.59)
is defined in $U$.
The space $\mathcal{L}(X; Y) =: Y_1$ is a normed space relative to which the mapping
(10.59) has the form (10.58), that is, $f' : U \to Y_1$, and it makes sense to speak
of differentiability for it.
If the mapping (10.59) is differentiable, its derivative
$$(f')' : U \to \mathcal{L}(X; Y_1) = \mathcal{L}(X; \mathcal{L}(X; Y))$$
is called the second derivative or second differential of $f$ and denoted $f''$ or $f^{(2)}$.
In general, we adopt the following inductive definition.
Definition 1 The derivative of order $n \in \mathbb{N}$ or $n$th differential of the mapping
(10.58) at the point $x \in U$ is the mapping tangent to the derivative of $f$ of order
$n-1$ at that point.
If the derivative of order $k \in \mathbb{N}$ at the point $x \in U$ is denoted $f^{(k)} (x)$, Definition 1
means that
$$f^{(n)} (x) := (f^{(n-1)})'(x).$$
Thus, if $f^{(n)} (x)$ is defined, then
$$f^{(n)}(x) \in \mathcal{L}(X; Y_n) = \mathcal{L}(X; \mathcal{L}(X; Y_{n-1})) =$$
$$= \cdots = \mathcal{L}(X; \mathcal{L}(X; \ldots; \mathcal{L}(X; Y)) \ldots).$$
(10.60)
Consequently, by Proposition 4 of Sect. 10.2, $f^{(n)} (x)$, the differential of order $n$ of
the mapping (10.58) at the point $x$ can be interpreted as an element of the space
$$\underbrace{\mathcal{L}(X, \ldots, X}_{n \text{ factors}}; Y)$$
of continuous $n$-linear transformations.
We note once again that the tangent mapping $f'(x) : T_Xx \to T_{Y_{f(x)}}Y$ is a map-
ping of tangent spaces, each of which, because of the affine or vector-space structure
of the spaces being mapped, we have identified with the corresponding vector space
and said on that basis that $f'(x) \in \mathcal{L}(X; Y)$. It is this device of regarding elements
$f'(x_1) \in \mathcal{L}(T_{x_1}X; T_{f(x_1)}Y)$ and $f'(x_2) \in \mathcal{L}(T_{x_2}X, T_{f(x_2)}Y)$, which lie in different
spaces, as vectors in the same space $\mathcal{L}(X; Y)$ that provides the basis for defining
higher-order differentials of mappings of normed vector spaces. In the case of an
affine or vector space there is a natural connection between vectors in the differ-
ent tangent spaces corresponding to different points of the original space. In the
final analysis, it is this connection that makes it possible to speak of the continuous
differentiability of both the mapping (10.58) and its higher-order differentials.